% BEGIN R28_MEMORY_PARITY
% Continuación aditiva de MEM01: los lectores ya están definidos.
\begin{proposition}[Paridad de los lectores bajo inversión orientada]
Con índices cíclicos en \(\mathbb Z/12\mathbb Z\), sea
\((C_M\tau)_m=-\tau_{-m}\) y \(\operatorname{sgn}(0)=0\).
Los dos lectores satisfacen
\[
 \nu_{120}(C_M\tau)=-\nu_{120}(\tau),\qquad
 \nu_{270}(C_M\tau)=\nu_{270}(\tau).
\]
\end{proposition}
\begin{proof}
La reflexión \(m\mapsto-m\) permuta las cuatro clases de posiciones
módulo cuatro. La inversión cambia el signo de cada suma de clase y,
por la imparidad de \(\operatorname{sgn}\), el de \(\nu_{120}\).
En el segundo lector se cancelan los signos de los dos factores:
\[
 \sum_m(C_M\tau)_m(C_M\tau)_{m+3}
 =\sum_m\tau_{-m}\tau_{-m-3}
 =\sum_k\tau_{k+3}\tau_k
 =\nu_{270}(\tau).
\]
\end{proof}
La orientación distingue así una lectura impar y una correlación par.
Ambas dependen de la disposición de las muestras en el registro; sus
leyes de transformación precisan qué información persiste al invertir
la ruta.
% END R28_MEMORY_PARITY
